
As promised, there are many different branches of discrete mathematics and each has its own flavor.  Each chapter of this textbook jumps discretely (that is, without a smooth transition) to a new topic.  As John Cleese from the TV show Monty Python's Flying Circus used to say when changing subjects ``and now for something completely different \ldots.'' 

This chapter introduces number theory.  Number theory is the backbone of modern encryption systems, the tools that keep information secret.  You can pay online securely with a credit card or look up your medical health information online because number theory is hard at work.  Anytime you visit a website that starts with \url{https}, the site is encrypted.  Number theory is also essential in building error-correcting codes, which insure that transmitted information is correctly received, and in powerful tools such as random number generators.

What is number theory?  It is the study of integers and their arithmetic properties.  We start in \Cref{sec_even_odd_divides} {\secintegerdivision} and \Cref{sec_div_alg} {\secdivalg} by discussing division in the integers.  Division is a more complicated operation than addition, subtraction, and multiplication because when you divide one integer by another, you might get a fraction that is not an integer. The remainders we get when dividing one integer by another form the basis for a new modular arithmetic we present in \Cref{sec_mod_arith} {\secmodarith}.  We close our exploration of number theory by diving deeper into the concept of divisors, including integers that have the minimum number of divisors in \Cref{sec_primes} {\secprimes} and divisors that two integers have in common in \Cref{sec_gcd} {\secgcd}.


